\booksubheading{Local Variable ও Local Scope}

\begin{definitionbox}

\textbf{Local Variable:}
কোনো Function-এর ভিতরে তৈরি এবং সাধারণত সেই Function-এর
মধ্যেই ব্যবহারযোগ্য Variable-কে \textbf{Local Variable} বলে।

\medskip

\textbf{Local Scope:}
Program-এর যে সীমিত অংশে একটি Local Variable ব্যবহার করা যায়,
তাকে সেই Variable-এর \textbf{Local Scope} বলে।

\end{definitionbox}

\begin{codeoutputbox}{Local Variable}

\noindent
\begin{minipage}[t]{0.52\linewidth}
\codeheading

\begin{english}
\ttfamily
def show():\\
\hspace*{4mm}x = 10\\
\hspace*{4mm}print(x)\\[0.5mm]
show()
\end{english}

\end{minipage}
\hfill
\begin{minipage}[t]{0.44\linewidth}
\outputheading

\begin{center}
\textenglish{\ttfamily 10}
\end{center}

\textenglish{\ttfamily x}
Function-এর ভিতরে তৈরি হয়েছে।

তাই
\textenglish{\ttfamily x}
একটি \textbf{Local Variable}।

\end{minipage}

\end{codeoutputbox}

%================================================
% Local outside
%================================================

\begin{codeoutputbox}{Local Variable বাইরে ব্যবহার করলে}

\noindent
\begin{minipage}[t]{0.52\linewidth}
\codeheading

\begin{english}
\ttfamily
def show():\\
\hspace*{4mm}x = 10\\[0.5mm]
show()\\
print(x)
\end{english}

\end{minipage}
\hfill
\begin{minipage}[t]{0.44\linewidth}
\outputheading

এখানে Error হবে।

কারণ
\textenglish{\ttfamily x}
শুধু
\textenglish{\ttfamily show()}
Function-এর Local Scope-এ তৈরি হয়েছে।

Function-এর বাইরে
\textenglish{\ttfamily x}
পরিচিত নয়।

\end{minipage}

\end{codeoutputbox}

\begin{rememberbox}

\begin{center}

\textbf{Function-এর ভিতরে তৈরি Variable}
\(\longrightarrow\)
\textbf{Local Variable}

\end{center}

সাধারণত সেটি সংশ্লিষ্ট Function-এর ভিতরেই ব্যবহারযোগ্য।

\end{rememberbox}

%================================================
% Parameter scope
%================================================

\booksubheading{Parameter-এর Scope}

\begin{definitionbox}

\textbf{Parameter Scope:}
Function Definition-এ ব্যবহৃত Parameter সাধারণত সেই Function-এর
ভিতরেই Local Name হিসেবে ব্যবহার করা যায়। Parameter-এর এই
ব্যবহারযোগ্য সীমাকে তার \textbf{Parameter Scope} বলা যায়।

\end{definitionbox}

\begin{codeoutputbox}{Parameter ও Local Variable}

\noindent
\begin{minipage}[t]{0.53\linewidth}
\codeheading

\begin{english}
\ttfamily
def square(n):\\
\hspace*{4mm}result = n * n\\
\hspace*{4mm}return result\\[0.5mm]
print(square(4))
\end{english}

\end{minipage}
\hfill
\begin{minipage}[t]{0.43\linewidth}
\outputheading

\begin{center}
\textenglish{\ttfamily 16}
\end{center}

\textenglish{\ttfamily n}
\(\longrightarrow\)
Parameter

\medskip

\textenglish{\ttfamily result}
\(\longrightarrow\)
Local Variable

\medskip

দুটিই
\textenglish{\ttfamily square()}
Function-এর Local Scope-এর অংশ।

\end{minipage}

\end{codeoutputbox}

%================================================
% Global
%================================================

\booksubheading{Global Variable ও Global Scope}

\begin{definitionbox}

\textbf{Global Variable:}
Function-এর বাইরে তৈরি এবং Program-এর বিভিন্ন অংশ থেকে
ব্যবহারযোগ্য Variable-কে সাধারণভাবে \textbf{Global Variable} বলে।

\medskip

\textbf{Global Scope:}
Program-এর যে বিস্তৃত অংশে একটি Global Name বা Variable
Access করা যায়, তাকে তার \textbf{Global Scope} বলে।

\end{definitionbox}

\begin{codeoutputbox}{Global Variable}

\noindent
\begin{minipage}[t]{0.52\linewidth}
\codeheading

\begin{english}
\ttfamily
x = 10\\[0.5mm]
def show():\\
\hspace*{4mm}print(x)\\[0.5mm]
show()
\end{english}

\end{minipage}
\hfill
\begin{minipage}[t]{0.44\linewidth}
\outputheading

\begin{center}
\textenglish{\ttfamily 10}
\end{center}

\textenglish{\ttfamily x}
Function-এর বাইরে Define করা হয়েছে।

তাই এটি \textbf{Global Variable}।

Function-এর ভিতর থেকে এর Value পড়া গেছে।

\end{minipage}

\end{codeoutputbox}

%================================================
% Local global comparison
%================================================

\booksubheading{Local বনাম Global Variable}

\begin{tabularx}
{\linewidth}
{
|>{\bfseries\RaggedRight\arraybackslash}p{0.28\linewidth}
|Y
|Y|
}

\hline
\rowcolor{softBlue}

\textbf{বিষয়}
&
\textbf{Local Variable}
&
\textbf{Global Variable}
\\
\hline

কোথায় তৈরি
&
Function-এর ভিতরে
&
Function-এর বাইরে
\\
\hline

Scope
&
সাধারণত সংশ্লিষ্ট Function
&
Program-এর বিস্তৃত অংশ
\\
\hline

প্রধান কাজ
&
Function-এর নিজস্ব Data
&
একাধিক অংশে প্রয়োজনীয় Data
\\
\hline

পরিবর্তনের প্রভাব
&
সাধারণত সীমিত
&
বিভিন্ন অংশে প্রভাব ফেলতে পারে
\\
\hline

\end{tabularx}

\begin{rememberbox}

\begin{center}

\textbf{Inside Function}
\(\longrightarrow\)
\textbf{Local}

\qquad

\textbf{Outside Function}
\(\longrightarrow\)
\textbf{Global}

\end{center}

\end{rememberbox}

%================================================
% Same name
%================================================

\booksubheading{একই নামের Local ও Global Variable}

একই Name-এর Local ও Global Variable থাকতে পারে।

\begin{codeoutputbox}{একই Name, ভিন্ন Scope}

\noindent
\begin{minipage}[t]{0.52\linewidth}
\codeheading

\begin{english}
\ttfamily
x = 10\\[0.5mm]
def show():\\
\hspace*{4mm}x = 20\\
\hspace*{4mm}print(x)\\[0.5mm]
show()\\
print(x)
\end{english}

\end{minipage}
\hfill
\begin{minipage}[t]{0.44\linewidth}
\outputheading

\begin{english}
\ttfamily
20\\
10
\end{english}

\medskip

Function-এর ভিতরের
\textenglish{\ttfamily x = 20}
হলো Local।

বাইরের
\textenglish{\ttfamily x = 10}
হলো Global।

দুটির Scope আলাদা।

\end{minipage}

\end{codeoutputbox}

\begin{rememberbox}

\textbf{\color{deepOrange}MCQ Trap:}

Function-এর ভিতরে একই Name-এর Local Variable থাকলে
Function-এর ভিতরে সাধারণত Local Variable-ই ব্যবহৃত হয়।

\end{rememberbox}

%================================================
% Read global
%================================================

\booksubheading{Global Variable পড়া}

Global Variable-এর Value Function-এর ভিতর থেকে পড়ার জন্য
সবসময়
\textenglish{\ttfamily global}
Keyword প্রয়োজন হয় না।

\begin{codeoutputbox}{Global Value Read}

\noindent
\begin{minipage}[t]{0.52\linewidth}
\codeheading

\begin{english}
\ttfamily
x = 10\\[0.5mm]
def show():\\
\hspace*{4mm}print(x)\\[0.5mm]
show()
\end{english}

\end{minipage}
\hfill
\begin{minipage}[t]{0.44\linewidth}
\outputheading

\begin{center}
\textenglish{\ttfamily 10}
\end{center}

এখানে শুধু Global
\textenglish{\ttfamily x}
-এর Value পড়া হয়েছে।

তাই
\textenglish{\ttfamily global x}
লেখার প্রয়োজন হয়নি।

\end{minipage}

\end{codeoutputbox}

%================================================
% global keyword
%================================================

\booksubheading{\textenglish{\ttfamily global} Keyword}

\begin{definitionbox}

\textbf{\textenglish{\ttfamily global} Keyword:}
Function-এর ভিতরে কোনো Name-কে Local Variable হিসেবে না ধরে
বাইরের Global Variable হিসেবে ব্যবহার ও Assignment করার ঘোষণা দিতে
\textenglish{\ttfamily global}
Keyword ব্যবহৃত হয়।

\end{definitionbox}

\begin{codeoutputbox}{\textenglish{\ttfamily global} ব্যবহার}

\noindent
\begin{minipage}[t]{0.52\linewidth}
\codeheading

\begin{english}
\ttfamily
x = 10\\[0.5mm]
def change():\\
\hspace*{4mm}global x\\
\hspace*{4mm}x = 20\\[0.5mm]
change()\\
print(x)
\end{english}

\end{minipage}
\hfill
\begin{minipage}[t]{0.44\linewidth}
\outputheading

\begin{center}
\textenglish{\ttfamily 20}
\end{center}

\textenglish{\ttfamily global x}
দিয়ে বোঝানো হয়েছে যে Function-এর
\textenglish{\ttfamily x}
নতুন Local Variable নয়।

Global
\textenglish{\ttfamily x}
-ই পরিবর্তিত হয়েছে।

\end{minipage}

\end{codeoutputbox}

%================================================
% Without global
%================================================

\booksubheading{\textenglish{\ttfamily global} ব্যবহার না করলে}

\begin{codeoutputbox}{Local Assignment}

\noindent
\begin{minipage}[t]{0.52\linewidth}
\codeheading

\begin{english}
\ttfamily
x = 10\\[0.5mm]
def change():\\
\hspace*{4mm}x = 20\\
\hspace*{4mm}print(x)\\[0.5mm]
change()\\
print(x)
\end{english}

\end{minipage}
\hfill
\begin{minipage}[t]{0.44\linewidth}
\outputheading

\begin{english}
\ttfamily
20\\
10
\end{english}

\medskip

Function-এর ভিতরের
\textenglish{\ttfamily x = 20}
একটি Local Variable তৈরি করেছে।

Global
\textenglish{\ttfamily x = 10}
অপরিবর্তিত আছে।

\end{minipage}

\end{codeoutputbox}

%================================================
% Good design
%================================================

\booksubheading{ভালো Function Design}

প্রয়োজন ছাড়া অনেক Global Variable-এর উপর নির্ভর না করে
Parameter দিয়ে Data নেওয়া এবং
\textenglish{\ttfamily return}
দিয়ে Result ফেরত দেওয়া সাধারণত বেশি পরিষ্কার Design।

\begin{codeoutputbox}{Parameter + Return}

\noindent
\begin{minipage}[t]{0.52\linewidth}
\codeheading

\begin{english}
\ttfamily
def add(a, b):\\
\hspace*{4mm}return a + b\\[0.5mm]
result = add(10, 20)\\
print(result)
\end{english}

\end{minipage}
\hfill
\begin{minipage}[t]{0.44\linewidth}
\outputheading

\begin{center}
\textenglish{\ttfamily 30}
\end{center}

Data এসেছে Parameter দিয়ে।

Result বাইরে গেছে
\textenglish{\ttfamily return}
দিয়ে।

এতে Function Global Variable-এর উপর কম নির্ভরশীল হয়।

\end{minipage}

\end{codeoutputbox}

\begin{rememberbox}

\begin{center}

\textbf{Parameter দিয়ে Input}
\(\longrightarrow\)
\textbf{Processing}
\(\longrightarrow\)
\textenglish{\ttfamily return}
\textbf{ দিয়ে Result}

\end{center}

এটি Function-কে আরও পরিষ্কার ও Reusable করতে সাহায্য করে।

\end{rememberbox}

%================================================
% Scope importance
%================================================

\booksubheading{Scope কেন গুরুত্বপূর্ণ?}

\begin{compacttwocolpoints}

\item Variable কোথায় ব্যবহার করা যাবে তা পরিষ্কার করে।

\item একই Name-এর Variable গুলিয়ে ফেলার সম্ভাবনা কমায়।

\item অপ্রয়োজনীয় Data Change কমায়।

\item Function-কে স্বাধীনভাবে কাজ করতে সাহায্য করে।

\item Program বুঝতে ও Debug করতে সহজ করে।

\end{compacttwocolpoints}

%================================================
% File
%================================================

